\documentclass[11pt,a4paper]{article}
\usepackage[utf8]{inputenc}
\usepackage[T1]{fontenc}
\usepackage[french]{babel}
\usepackage{geometry}
\usepackage{booktabs}
\usepackage{amsmath}
\usepackage{enumitem}
\geometry{margin=2.5cm}

\title{Calibration du vote majoritaire du garde-fou de transmission\\
	\large Note méthodologique pour les rédacteurs de l'article}
\author{Projet AgriSense}
\date{}

\begin{document}
	\maketitle
	
	\section{Objet de la note}
	Cette note décrit comment les seuils du garde-fou ont été calibrés, quels résultats ont été obtenus et
	quels seuils sont retenus. Le mécanisme n'est pas modifié : une image est transmise si \textbf{au moins
		deux critères sur trois} dépassent leur seuil (vote majoritaire). Seule la valeur des seuils est optimisée.
	
	\section{Le mécanisme}
	Pour chaque paire (image précédente, nouvelle image), trois critères sont calculés.
	\begin{itemize}[nosep]
		\item \textbf{HIST} : distance de Bhattacharyya entre les histogrammes couleur globaux (4~classes par canal RGB).
		\item \textbf{MEAN} : écart moyen, sur les trois canaux, entre les médianes des deux images.
		\item \textbf{PROP} : pourcentage de blocs de $32\times32$ pixels dont la distance de Bhattacharyya
		dépasse un \emph{seuil de bloc}. Une image $640\times480$ contient 300 blocs.
	\end{itemize}
	Chaque critère vote « transmettre » lorsqu'il est \textbf{strictement supérieur} à son seuil. PROP comporte
	donc deux paramètres : le seuil de bloc (décision pour chaque bloc) et le seuil PROP (décision pour l'image
	entière). Au total, quatre valeurs sont calibrées : HIST, MEAN, PROP et le seuil de bloc.
	
	\section{Données}
	Le jeu de calibration contient 2\,503 paires d'images étiquetées :
	\begin{center}
		\begin{tabular}{lrl}
			\toprule
			Étiquette & Paires & Signification \\
			\midrule
			0 & 840 & auto-paires (image comparée à elle-même), métriques nulles \\
			1 & 831 & images similaires (ne pas transmettre) \\
			2 & 832 & images différentes (transmettre) \\
			\bottomrule
		\end{tabular}
	\end{center}
	La classe positive est l'étiquette~2. Deux scénarios sont rapportés :
	\begin{itemize}[nosep]
		\item \textbf{hard} (scénario principal) : étiquettes 1 et 2 uniquement, $n=1\,663$. Les auto-paires sont exclues
		car leurs métriques valent 0 et rendent le problème trivial.
		\item \textbf{all} : toutes les étiquettes, $n=2\,503$. Il est optimiste et n'est donné qu'à titre indicatif.
	\end{itemize}
	Le jeu est équilibré par construction : il ne reflète pas la fréquence réelle des changements sur le terrain.
	Aucune estimation de cette fréquence n'a été faite.
	
	\section{Méthode de calibration}
	
	\subsection{Recherche des seuils}
	Pour chaque critère, une grille de 25 valeurs est construite à partir des percentiles 2 à 98 des données
	d'entraînement. Toutes les combinaisons (HIST, MEAN, PROP) sont évaluées pour chacun des cinq seuils de bloc
	candidats (0,05 ; 0,1 ; 0,2 ; 0,3 ; 0,5). Le point retenu est le centre du plateau quasi-optimal (tolérance
	0,002), plutôt que le premier maximum, afin de favoriser la stabilité.
	
	\subsection{Objectifs opérationnels}
	\begin{itemize}[nosep]
		\item \texttt{bacc} : maximiser l'exactitude équilibrée, moyenne du recall (vrais changements détectés) et de la
		spécificité (images sans changement écartées).
		\item \texttt{spec@0.95} et \texttt{spec@0.98} : maximiser la spécificité sous une contrainte de recall d'au moins 0,95
		ou 0,98.
	\end{itemize}
	
	\subsection{Deux méthodes de calibration comparées}
	\begin{description}[nosep]
		\item[Classique.] Le meilleur point de grille est choisi sur tout l'entraînement. Pour les objectifs \texttt{spec@}, la
		contrainte de recall est testée sur la borne basse de Wilson ($z=1{,}645$). Le seuil de bloc est choisi par
		validation croisée interne (3~folds).
		\item[Robuste.] Chaque point de grille est évalué sur 200 rééchantillons de 20 blocs temporels de l'entraînement.
		On retient le point dont la performance est atteinte dans 90\,\% des rééchantillons, et non seulement en moyenne. Le
		seuil de bloc est choisi dans la même procédure.
	\end{description}
	
	\subsection{Validation}
	La performance est estimée hors échantillon par \textbf{validation croisée imbriquée} :
	\begin{itemize}[nosep]
		\item cinq folds externes découpés en \textbf{blocs temporels} (ordre des numéros d'image) ;
		\item \textbf{purge} : toute paire d'entraînement partageant une image avec le fold de test est retirée
		(entre 135 et 224 paires selon le fold), car les paires forment une composante connexe et un découpage aléatoire
		laisserait fuir de l'information ;
		\item les seuils de chaque fold sont calculés sans jamais voir ses données de test.
	\end{itemize}
	Les intervalles de confiance à 95\,\% sont obtenus par bootstrap par blocs temporels (2\,000 rééchantillons, 40 blocs).
	Les différences entre méthodes sont appariées sur les mêmes rééchantillons.
	
	\section{Résultats}
	
	\subsection{Pouvoir discriminant des critères (scénario \textit{hard})}
	\begin{center}
		\begin{tabular}{lc}
			\toprule
			Critère & AUC [IC 95\,\%] \\
			\midrule
			HIST & 0,958 [0,947 ; 0,970] \\
			MEAN & 0,876 [0,856 ; 0,897] \\
			PROP (bloc 0,05) & 0,987 [0,981 ; 0,992] \\
			PROP (bloc 0,1) & 0,987 [0,980 ; 0,992] \\
			\bottomrule
		\end{tabular}
	\end{center}
	PROP est le critère le plus discriminant, MEAN le moins.
	
	\subsection{Performances hors échantillon (scénario \textit{hard})}
	\begin{center}
		\small
		\begin{tabular}{llcccc}
			\toprule
			Objectif & Méthode & Recall & Spécificité & Bal.\ acc. & Recall du pire fold \\
			\midrule
			\texttt{bacc} & classique & 0,944 & 0,961 & 0,953 & 0,871 \\
			\texttt{bacc} & robuste   & 0,944 & 0,954 & 0,949 & 0,806 \\
			\addlinespace
			\texttt{spec@0.95} & classique & 0,939 & 0,929 & 0,934 & 0,742 \\
			\texttt{spec@0.95} & robuste   & 0,936 & 0,948 & 0,942 & 0,742 \\
			\addlinespace
			\texttt{spec@0.98} & classique & 0,987 & 0,798 & 0,892 & 0,944 \\
			\texttt{spec@0.98} & robuste   & 0,982 & 0,845 & 0,913 & 0,935 \\
			\bottomrule
		\end{tabular}
	\end{center}
	Intervalles de confiance à 95\,\% pour \texttt{bacc} classique : recall $[0{,}917 ; 0{,}966]$, spécificité $[0{,}940 ; 0{,}980]$.
	
	Principales observations :
	\begin{itemize}[nosep]
		\item En \texttt{bacc}, la méthode robuste ne fait pas mieux que la classique (différences non significatives) et son
		pire fold est moins bon.
		\item Aux objectifs \texttt{spec@}, la méthode robuste gagne en spécificité à recall comparable :
		$+0{,}019$ $[+0{,}003 ; +0{,}042]$ pour \texttt{spec@0.95} et $+0{,}047$ $[+0{,}028 ; +0{,}068]$
		pour \texttt{spec@0.98} (intervalles excluant 0).
		\item La cible de recall de 0,95 n'est tenue en moyenne par aucune méthode (0,936 à 0,939), et seuls 3~folds sur~5
		l'atteignent. Ce réglage ne doit pas être présenté comme garantissant 95\,\%.
	\end{itemize}
	
	\subsection{Erreurs réelles du réglage retenu}
	Décisions hors échantillon de la méthode classique, objectif \texttt{bacc}, scénario \textit{hard} :
	\begin{center}
		\begin{tabular}{lcc}
			\toprule
			& Bloquées (0) & Envoyées (1) \\
			\midrule
			Paires similaires (831) & 799 (96,1\,\%) & \textbf{32} (3,9\,\%) \\
			Paires différentes (832) & \textbf{47} (5,6\,\%) & 785 (94,3\,\%) \\
			\bottomrule
		\end{tabular}
	\end{center}
	Les 32 paires similaires envoyées sont des transmissions inutiles. Les 47 paires différentes bloquées sont des
	changements non transmis, l'erreur la plus grave pour l'application. Avec la méthode robuste, on obtient 38 et 47.
	
	\section{Seuils retenus}
	Le réglage retenu est celui de l'objectif \texttt{bacc}, calibré par la méthode classique sur l'ensemble du scénario
	\textit{hard} :
	\begin{center}
		\begin{tabular}{lc}
			\toprule
			Paramètre & Valeur \\
			\midrule
			Seuil HIST & $> 0{,}00275$ \\
			Seuil MEAN & $> 16{,}67$ \\
			Seuil PROP & $> 64\,\%$ des blocs \\
			Seuil de bloc (dans PROP) & $0{,}05$ \\
			Règle de décision & au moins 2 critères sur 3 \\
			\bottomrule
		\end{tabular}
	\end{center}
	Avec 300 blocs, un seuil de 64\,\% signifie plus de 192 blocs, soit au moins 193 blocs « changés ».
	
	\paragraph{Justification.} Ce réglage a le meilleur compromis global hors échantillon en \texttt{bacc}, et le pire fold le
	moins mauvais. Le seuil de bloc 0,05 est sélectionné dans les 5~folds, et le seuil PROP varie peu (55,5 à 69,9\,\%
	selon les folds). L'objectif \texttt{bacc} traite de façon symétrique les transmissions inutiles et les changements
	non transmis.
	
	\paragraph{Alternative prudente.} Pour ne presque rien rater, l'objectif \texttt{spec@0.98} détecte environ 98,7\,\% des vrais
	changements (classique) au prix de 20\,\% de transmissions inutiles en plus. Les seuils de la méthode robuste étant
	instables pour cet objectif (intervalles de confiance très larges sur HIST et MEAN), ce réglage doit être présenté avec réserve.
	
	\section{Limites à mentionner dans l'article}
	\begin{itemize}
		\item \textbf{Les gains ne sont pas mesurés en énergie.} Aucune mesure énergétique n'a été faite : l'article doit
		exprimer les bénéfices en \emph{transmissions inutiles évitées} (environ 96\,\% des paires similaires bloquées), non en énergie.
		\item \textbf{Fréquence réelle des changements inconnue.} Le jeu est équilibré : le taux global de transmissions
		évitées sur le terrain n'est pas estimable avec ces données.
		\item \textbf{Variabilité entre périodes.} Le pire fold rate 12,9\,\% des vrais changements (recall 0,871), contre 5,6\,\%
		en moyenne.
		\item \textbf{Rôle limité de MEAN.} Dans le réglage retenu, MEAN n'est décisif que dans environ 12\,\% des paires : le vote
		se joue surtout entre HIST et PROP.
		\item \textbf{Scénario \textit{all}.} Il inclut les auto-paires triviales et surestime les performances ; il ne doit pas
		être utilisé comme résultat principal.
		\item \textbf{Comparaison avec l'existant non réalisée.} L'amélioration par rapport aux seuils actuels de l'article n'a pas encore
		été mesurée (voir ci-dessous).
	\end{itemize}
	
	\section{À faire avant soumission}
	\begin{enumerate}[nosep]
		\item Relancer la calibration avec \texttt{--article-thresholds HIST,MEAN,PROP --article-block B} pour chiffrer le gain par
		rapport aux seuils actuels, avec un intervalle de confiance.
		\item Vérifier que le firmware reproduit exactement les métriques du script : résolution $640\times480$, blocs $32\times32$,
		histogrammes à 4 classes par canal, distance de Bhattacharyya, comparaison stricte $>$.
		\item Idéalement, valider le réglage retenu sur un jeu collecté séparément (autre période ou autre site).
	\end{enumerate}
	
\end{document}